\section{Discussion}

\paragraph{Forced-choice grounding as a GUI search failure mode.}
The central finding is that target-present GUI search models can behave like forced-choice grounding systems when the target is absent. They may generate short, confident, and convergent paths even when the requested target is not visible. This is different from ordinary localization error: the model is not merely off by some distance, but is making a target-presence decision without being evaluated for it.

\paragraph{Why path and OCR evidence help.}
Path evidence captures whether the generated search behavior actually approaches a target-like region. OCR evidence captures whether visible UI text supports the target cue. Neither signal is sufficient alone: path evidence can be brittle when the generated scanpath is short or off-target, and OCR can miss small or stylized text. The combined AND verifier works because it rejects only when both signals are weak, making it an interpretable safety layer rather than a replacement for the base model.

\paragraph{How to interpret the VLM results.}
Generic VLM yes/no prompting is a strong baseline and should remain in the paper. It shows that target presence can be partly solved by screenshot reasoning alone. However, prompt variants differ substantially, which means VLM presence classification is itself a modeling choice. Evidence-aware VLM verification is stronger because it combines screenshot reasoning with structured path/OCR evidence and now holds up on the matched image split. The paper should therefore distinguish two claims: combined AND is the strongest interpretable verifier, while evidence-aware VLM is the strongest validated verifier so far.

\paragraph{Implications for non-text and associative targets.}
The released subset is text-target dominated, so the current results do not solve non-text, icon, or associative GUI search. However, the target-absence formulation is relevant to those directions. Non-text target search will need candidate evidence beyond OCR, such as icon proposals, visual embeddings, or crop-level VLM verification. Associative search will need new query annotations because correctness is more semantic and context-dependent.
